\documentclass[11pt]{article}
\usepackage[letterpaper,margin=.86in]{geometry}
\usepackage{amsmath,amssymb,amsthm,hyperref}
\hypersetup{hidelinks,pdftitle={The complete zero sum at the quadratic nonresidue scale}}
\newcommand{\R}{\mathbb R}
\newcommand{\e}{\mathrm e}
\newcommand{\ii}{\mathrm i}
\newcommand{\dd}{\,\mathrm d}
\newcommand{\ind}{\mathbf1}
\newtheorem{proposition}{Proposition}
\title{The complete zero sum at the quadratic nonresidue scale}
\author{Research notes}
\date{}
\begin{document}
\maketitle
Assume Dirichlet GRH, including RH for $\zeta$. Let $p$ be an odd prime,
$\chi=(\cdot/p)$, $L=\log p$, $y=cL^2$ and $U=\sqrt L$, with $c>0$
fixed. The goal is $n_p=o(L^2)$. Ordinates of the nontrivial zeros
of $L(s,\chi)$ are counted with multiplicity. Use the two statistics
\begin{align*}
 F^+(y,U)&=\int_{\R}\left|\sum_{0\le\gamma\le U}
               \e^{\ii(\log y+v)\gamma}\right|^2\e^{-2|v|}\dd v,\\
 Z(t)&=\sum_{\gamma}\frac{2y^{\ii\gamma}}{1+(t-\gamma)^2},\qquad
 Q_p(c)=\frac{1}{\pi y}\int_0^U|Z(t)|^2\dd t .
\end{align*}
These normalizations differ. We prove $F^+(y,U)\ll_c L^2$ and the
exceptional lower coefficient $1/12$ for $F^+/y$. For the complete
statistic, $Q_p(c)\le 4/3+o_c(1)$ for every prime. On an exceptional
sequence the portion of its norm forced by the primes below $y$ is
$1/3+o_c(1)$. The cross term with the unknown tail tends to zero.

\section{The negative-shift test for the sharp pair sum}
In this paragraph only, $U=U(p)$ can be any height tending to infinity.
Let $g$ be real and smooth with compact support in $(-\infty,0)$ and
put $\Phi(t)=\int g(v)\e^{\ii tv}\dd v$. The explicit formula for
a primitive real character \cite[Section 3]{cmqr} gives, if $n_p>y$,
\[
 \frac{1}{\sqrt y}\sum_\gamma y^{\ii\gamma}\Phi(\gamma)
 =-\frac{1}{\sqrt y}\sum_n \frac{\Lambda(n)}{\sqrt n}
                 g(\log(n/y))+o_g(1)
 \longrightarrow-\int_{-\infty}^0 g(v)\e^{v/2}\dd v .
\]
The conductor and reflected prime terms vanish for large $p$, the
gamma integral is $O_g(1)$ before normalization, and the prime number
theorem gives the limit. The tail past $|\gamma|>U$ before normalization
is $O_{g,N}((L+\log(2+U))(1+U)^{-N})$ by the unit-interval zero count.
The central multiplicity is $O(L/\log L)$ by (3.3)--(3.4) of \cite{cmqr},
using the limit of the half-weight count at zero. Adding the positive
sum to its conjugate counts those central entries twice, an $o_c(\sqrt y)$
correction. Weighted Cauchy--Schwarz now gives
\[
 \liminf_{\rm exc}\frac{F^+(y,U)}y\ \ge\
 \frac14\sup_{g\ne0}
 \frac{|\int g(v)\e^{v/2}\dd v|^2}{\int |g(v)|^2\e^{2|v|}\dd v}
 =\frac14\int_{-\infty}^0\e^{3v}\dd v=\frac1{12}.
 \tag{1}
\]
The supremum follows by smooth compact approximation of
$\e^{5v/2}\ind_{v<0}$ in the denominator norm. Hence a uniform upper
coefficient $0.08$ for $F^+/y$ at each fixed dyadic $c$ would suffice.

\section{Grouping the prime powers at the chosen height}
For $n\ge2$ put $s_n=\log(n/y)$ and
\[
 b_n=\begin{cases}\Lambda(n)\sqrt n/y,&n\le y,\\
                    \Lambda(n)y/n^{3/2},&n>y.\end{cases}
 \qquad
 \mu_\chi=y^{-1/2}\sum_n b_n\chi(n)\delta_{s_n}.
\]
Write $\mu_1$ for the all-positive comparison, and use the angular
Fourier transform $\widehat\mu(t)=\int\e^{-\ii ts}\dd\mu(s)$.
Then $-\sqrt y\,\widehat\mu_\chi(t)$ is the prime term in (28) of
Kandhil, Languasco and Moree \cite{klm}. Its series converges absolutely.
Their identities (14)--(16), (24)--(25) bound the norm of the remaining
terms on $[0,U]$ by
$O(\sqrt U(L+\log(2+U))/y+1)$. Any half-weight convention at $n=y$
costs $O(\sqrt U\log(2y)/\sqrt y)$ in the same norm. Both errors
are $o_c(\sqrt y)$. The identities are valid here without a
pair-correlation hypothesis; the asymptotic theorem in \cite{klm}
with $T\ge x$ is not used.

Here is a logarithmic grouping argument that retains the leading
coefficient. Use a sliding interval
$k_h(s)=h^{-1}\ind_{[-h/2,h/2]}(s)$. With
$\operatorname{sinc}(t)=\sin(t)/t$, its Fourier multiplier is
$\operatorname{sinc}(th/2)$. The limiting unsigned density in
logarithmic position is
\[
 f(s)=\begin{cases}\e^{3s/2},&s\le0,\\ \e^{-s/2},&s>0.\end{cases}
 \qquad \|f\|_2^2=1/3+1=4/3 .
\]
For either $h=1/U$ or $h=1/(U\log y)$ one has
\[
 \|\mu_1*k_h-f\|_2=o_c(1),\qquad
 |\mu_\chi*k_h|\le \mu_1*k_h .
 \tag{2}
\]
To check the first claim, put $x=y\e^s$. For $x\ge y^{3/4}$ the
window is $[x\e^{-h/2},x\e^{h/2}]$. The estimate
$\psi(x)=x+O(\sqrt x\log^2(2x))$ gives mass $hx$ with relative error
$O(h+\log^2(2x)/(h\sqrt x))=o_c(1)$ uniformly in this range.
The factor $b_n/\Lambda(n)$ on this window varies relatively by
$O(h)$, including across $y$. Hence the density is $f(s)(1+o_c(1))$.
For the remaining $s$, the total normalized mass of indices that can
occur is $O(y^{-3/8})$ by Chebyshev. Young's inequality bounds their
smoothed $L^2$ norm by $O(y^{-3/8}h^{-1/2})=o(1)$, and $f$ has
vanishing $L^2$ norm there as well. The same proof applies to the
unsigned measure restricted to $s_n\le d$ or $s_n>d$, for any fixed
$d$. The boundary strip of width $h$ has vanishing squared norm.

With $h=1/U$, the multiplier is bounded away from zero for $|t|\le U$.
Plancherel and (2) therefore give $\int_0^U|\widehat\mu_\chi|^2=O_c(1)$.
This proves the stated order bound for the complete sum using precisely
windows of width $1/U$. For the coefficient choose $h=1/(U\log y)$, so
$Uh=o_c(1)$. The multiplier is then $1+o_c(1)$ throughout the band.
All the measures are real, so the two half-bands have equal energy.
If $\Pi_U$ denotes orthogonal Fourier projection onto $[-U,U]$, then
\[
 Q_p(c)=\|\Pi_U(\mu_\chi*k_h)\|_2^2+o_c(1)
 \ \le\ 4/3+o_c(1).
 \tag{3}
\]
The $o_c(1)$ is uniform in the prime. Applying the multiplier and
projection calculation to $\mu_1$ gives
$\pi^{-1}\int_0^U|\widehat\mu_1|^2\dd t\to4/3$ by (2).
This comparison uses only the prime coefficients.

\paragraph{The endpoints for the sharp statistic.}
Let $Z_I$ denote the sum defining $Z$ restricted to $\gamma\in I=[0,U]$.
The elementary convolution of the two Cauchy kernels gives
$\|Z_I\|_{L^2(\R)}^2=2\pi F^+(y,U)$. The unit-interval zero count gives
\[
 \|Z_I-\ind_I Z\|_{L^2(\R)}\ll L+\log(2+U).
 \tag{4}
\]
Inside $I$, bound the two omitted tails in absolute value by a constant
times $(L+\log(2+U))((1+t)^{-1}+(1+U-t)^{-1})$. For instance a
tail at distance $d$ is bounded by summing
$(L+\log(2+U+j))/(1+(d+j)^2)$ over integers $j\ge0$; on $I$ one
has $d\le U$. Outside $I$ the included ordinates contribute at most
$O((L+\log(2+U))/(1+\operatorname{dist}(t,I)))$. All these profiles
have bounded squared integrals. The triangle inequality at the norm
level followed by (3) gives $F^+(y,\sqrt L)\ll(\sqrt y+L+\log(2+U))^2
\ll_c L^2$. The endpoint error in (4) is not used in (3) or below.

\section{Separating the determined and undetermined signs}
Allow a fixed $A\ge1$ and suppose $n_p>Ay$. Split the measure
$\mu_\chi=\mu_<+\mu_>$ at $s=d=\log A$, assigning the boundary to
the lower part. Let
\[
 B_{\chi,A}=\frac1\pi\int_0^U|\widehat{\mu_>}(t)|^2\dd t .
\]
On any such exceptional sequence,
\[
 \boxed{\quad Q_p(c)=\frac43-\frac1A+B_{\chi,A}+o_{c,A}(1),
 \qquad 0\le B_{\chi,A}\le\frac1A+o_{c,A}(1).\quad}
 \tag{5}
\]
Indeed $\mu_<$ is unsigned, so $\Pi_U(\mu_<*k_h)\to
f\ind_{s\le d}$ strongly in $L^2$ by (2) and $U\to\infty$. The tail
$\mu_>*k_h$ has bounded norm, is supported in $[d-h/2,\infty)$,
and is dominated by the unsigned smoothed tail. Its inner product
with $f\ind_{s\le d}$ tends to zero. Self-adjointness and idempotence
of $\Pi_U$ thus make the real cross term of the two band-limited
parts $o_{c,A}(1)$. The upper tail norm squared is at most
$\int_d^\infty f^2+o(1)=1/A+o(1)$.
The lower norm squared tends to $\int_{-\infty}^d f^2=4/3-1/A$.
This proves (5) directly for the complete statistic in (16) of \cite{klm}.
In particular $A=1$ forces $\liminf_{\rm exc}Q_p(c)\ge1/3$.
For the raw integral the lower coefficient is $\pi/3$ times $y$.

\paragraph{The signed sum being bounded.}
Let $D_U(s)=\sin(Us)/s$ with $D_U(0)=U$. The exact prime-side energy is
\[
 \frac{1}{\pi y}\sum_{m,n} b_m b_n\chi(mn)D_U(\log(m/n)).
 \tag{6}
\]
This is the square in (3) before grouping and up to its $o_c(1)$
zero-formula error. The ordering is over every pair. Relative to
the all-positive comparison a pair of nonzero signs changes its
contribution precisely when one is negative and the other positive.
Under $n_p>y$ a negative prime-power sign requires an odd power of
a nonresidue prime exceeding $y$. Two such negative signs multiply
to $+1$. Terms with an index divisible by $p$ have zero sign.
Ramified prime powers have negligible mass, as does the prime $p$
itself since $p>y$. More generally, all proper prime powers together
have $\sum b_n=O(1)$, by the bound
$\sum_{\ell^k\le X,\,k\ge2}\log\ell=O(\sqrt X)$ and dyadic summation
above $y$. Dropping them changes the normalized band norm by
$O(\sqrt{U/y})=o(1)$. This justifies the prime-only sums that follow.

For the full range, let $S$ be the sum of $b_m b_nD_U(\log(m/n))$
over primes $m,n$ with $\chi(m)=+1$ and $\chi(n)=-1$, each opposite
pair in this order only. Then (3) and (6) give
\[
 Q_p(c)=\frac43-\frac{4S}{\pi y}+o_c(1).
 \tag{7}
\]
The factor four includes the sign flip and the reversed ordered pair.
A uniform upper coefficient $1/3-\epsilon_c$ therefore asks for
$4S/(\pi y)\ge1+\epsilon_c+o_c(1)$. The kernel in this sum changes
sign. On an exceptional sequence the part with $m\le y<n$ is $o_c(y)$:
the negative-prime measure above $y$ is half the difference of the
unsigned and signed tails, up to the negligible prime powers and
ramification, and the orthogonality argument proving (5) applies to both.
Writing $S_{>>}$ for the remaining sum with both primes above $y$, we have
\[
 \frac{4S_{>>}}{\pi y}=1-B_{\chi,1}+o_c(1),
 \qquad -o_c(1)\ \le\frac{4S_{>>}}{\pi y}\le1+o_c(1).
 \tag{8}
\]
Indeed the all-positive tail energy tends to $1$. Its difference
from $B_{\chi,1}$ changes exactly the opposite pairs in that tail.
Thus even complete cancellation of the high-prime energy leaves
the low-prime contribution $1/3$.

Both ends of the tail energy range are possible if only the magnitude
of each prime sign is used. All positive signs above $y$ give the
upper end. Independently assigned fair prime signs give expected
tail energy $U\sum_{\ell>y} b_\ell^2/(\pi y)\ll U\log y/y=o(1)$,
by Chebyshev on dyadic intervals. Proper powers cost $o(1)$ in norm
under any completely multiplicative extension. These examples test
the estimates; their signs are not asserted to come from primes $p$
of the required size.

\section{A smaller saving would suffice across the dyadic scales}
Suppose there is a fixed $\eta>0$ and infinitely many $p$ with
$n_p>\eta L^2$. For every fixed dyadic $c\le\eta$, apply (5) to the
same primes with $A=\eta/c$. Then
\[
 \liminf_{\rm exc} Q_p(c)\ \ge\ \frac43-\frac c\eta.
 \tag{9}
\]
Hence it suffices to prove, at each of a sequence of dyadic $c$ tending
to zero, an all-prime upper coefficient $4/3-\delta(c)$ with
$\delta(c)/c$ unbounded along that sequence. Even a fixed positive
gap along arbitrarily small dyadic scales would suffice. Each
limit in $p$ is taken first, so this uses no uniform error in $c$.

There is also a version with nonnegative pair weights. With the fine
width $h=1/(U\log y)$ put $J_h(s)=h^{-1}(1-|s|/h)_+$. The unsigned
smoothed density has squared norm $4/3+o_c(1)$. Subtracting the
squared norm of the signed smoothed density and using $k_h*k_h=J_h$
gives, by (3),
\[
 Q_p(c)\le\frac43-
  \frac4y\sum_{\substack{m,n\ {\rm primes}\\\chi(m)=1,\ \chi(n)=-1}}
       b_m b_n J_h(\log(m/n))+o_c(1).
 \tag{10}
\]
For this full-frequency density calculation, discarding proper powers
costs $O((hy)^{-1/2})$ in norm by Young's inequality. The prime $p$
also gives $o(1)$. Thus a weighted supply of nearby opposite-sign
primes could give the saving in (9).

\paragraph{How much mixing can GRH actually force this way?}
There is a positive contribution to (10) that can be proved for
the individual character. Fix a smooth $0\le g\le1$, nonzero and
supported in $(-1/2,1/2)$. The explicit formula gives, for $x>2$
and the above shrinking $h$,
\[
 \left|\sum_n\frac{\Lambda(n)\chi(n)}{\sqrt n}
              g\!\left(\frac{\log(n/x)}{h}\right)\right|
 \ll_g L+\log(2+1/h).
 \tag{11}
\]
In fact the zero test has absolute value
$O_g(h(1+h|\gamma|)^{-2})$, which sums to the right side by the
unit-interval count. The gamma integral is $O_g(\log(2+1/h))$;
the conductor and reflected prime terms vanish by support. The
unsigned prime sum with the same weights is
$h\sqrt x\int g(u)\e^{hu/2}\dd u+O_g(\log^2(2x))$ by RH and
partial summation. Proper powers remove $O_g(1)$ for $x$ in a
dyadic range. Take $X=D_g L^2/h^2$ with $D_g$ sufficiently large.
Uniformly for $X\le x\le2X$ the positive and negative primes each
contribute $\gg_g h\sqrt x$ in these weights, since
$\log^2(2X)=o(L)$ and $2X\e^h<p$ for large $p$.

Let $\mu_+,\mu_-$ be the normalized positive measures of the primes
of the respective signs, using $b_n/\sqrt y$ for both. For $x=y\e^s$
in that range the last estimates imply
$(\mu_\pm*k_h)(s)\gg_g\sqrt{y/x}$ since $b_n=y\Lambda(n)/n^{3/2}$
there. The whole subtraction in (10) is
$4\int(\mu_+*k_h)(\mu_-*k_h)\dd s$. Its portion on this dyadic
range is therefore $\gg_g y/X$; the unsigned prime estimate bounds
that same portion by $O_g(y/X)$. In particular it is of order
$c h^2=c/(U^2\log^2 y)$, which vanishes as $p$ grows at fixed $c$.
The factor $\log y$ in the fine width can be replaced throughout by
any more slowly diverging positive factor. The resulting contribution
from this range still vanishes.
This proves actual opposite signs on the log intervals, but at
scales carrying too little energy to yield the required coefficient.
Equations (8)--(11) give no fixed positive gap of the size in (9).
The bound for $n_p$ remains unchanged.

\begin{thebibliography}{9}
\bibitem{cmqr} E.~Carneiro, M.~Milinovich, E.~Quesada-Herrera and
A.~Ramos, \emph{Fourier optimization, the least quadratic non-residue,
and the least prime in an arithmetic progression},
\href{https://arxiv.org/html/2404.08380v2#S3}{Section 3}.
\bibitem{klm} N.~Kandhil, A.~Languasco and P.~Moree,
\emph{Beyond the Riemann Hypothesis bounds: A pair-correlation approach
to the least prime in arithmetic progression and the smallest quadratic
non-residue}, \href{https://arxiv.org/html/2607.14515v1#S3}{Sections 2--3}.
\end{thebibliography}
\end{document}
